\section{Introduction}
\label{sec:intro}
A question about a company, a product, and a news event can require several linked entities even when the final answer is a single name. In knowledge-based access to Web information, an intermediate selection step must retain these entities before an answer generator can use their relations. Selecting plausible endpoint names is insufficient when intermediate evidence is omitted; selecting many connected names can instead admit irrelevant neighbours. Understanding these errors requires examining the structure of the requested evidence, not only average selection accuracy.

We study this selection step on a fixed candidate pool drawn from a corporate knowledge graph (KG). The resource connects companies, people, products, locations, and news and was developed for corporate GraphRAG evaluation~\cite{kirpichev2026towards}. Our Web-related question is how reliably language models select the connected entity context needed to access and combine this structured information. We evaluate agreement with an explicitly defined entity-route reference, not the performance of a deployed Web search system or the correctness of generated answers.

Graph encodings and graph-pattern comprehension are established subjects of LLM evaluation~\cite{fatemi2024talklikegraph,dai2025graphpatterns}. The closely related Tool Graph Complexity manuscript already includes KG entity selection, motif and size comparisons, and adjacent distractors in Appendix~O~\cite{anonymousToolGraph}. We take that KG line as a related experimental foundation and examine structural profiles under distinguished evidence and presentation conditions. We do not claim that motif-based KG selection, size effects, or distractor sensitivity are new phenomena. The precise extent of shared historical task/response records with that appendix has not been independently established, and we make no independent-replication or new-resource claim.

Our analysis has three questions. \textbf{RQ1:} At comparable node counts, which graph-pattern task families are difficult for each model, and what selection errors occur? \textbf{RQ2:} Within a family, how does quality vary across available sizes and evidence conditions? \textbf{RQ3:} How robust are these profiles to explicit relation availability and adjacent distractors? Rendering, card placement, wording, and naming provide supplementary diagnostics rather than separate headline research questions.

The study makes three scoped contributions: (i) a descriptive characterization of pattern-family profiles, separating equal-size comparisons from families evaluated at different sizes; (ii) an all-model analysis of family-specific size profiles under native cards and their explicit relation overlay; and (iii) a reference-aware account of condition sensitivity, partial versus exact recovery, and omission versus inclusion errors. Audited pairing and failure-inclusive reporting make these comparisons inspectable. These are empirical contributions on a fixed KG, not a new selection algorithm.

The principal observations resist a single complexity ordering. At six nodes, outward stars are recovered much more often than inward trees, especially for Qwen. Increasing size is associated with substantial losses for some chains and trees but almost no loss for outward stars in some model/condition combinations. Explicit relations substantially increase recovery of the graph reference, while adjacent contexts remain difficult. Since motif families also differ in query semantics and accepted task composition, these results are conditional profiles rather than isolated causal effects of topology.
